% Converted from PROOF.md into LaTeX
\section{Proof of Non-Removability of the Seam Residue}

\subsection*{Object}
Finite nerve $N$ with regions $U_1, U_2, U_3, U_4$ and oriented edges
\[
U_1U_2,\quad U_2U_3,\quad U_3U_4,\quad U_1U_4.
\]

No 2-simplices: $C^2 = 0$.

Cochain complex over $\mathbb{Q}$:
\[
C^0 = \mathbb{Q}^4
\xrightarrow{\delta^0}
C^1 = \mathbb{Q}^4
\xrightarrow{\delta^1 = 0}
0.
\]

Residue:
\[
r = (1,\,1,\,1,\,-2) \in C^1.
\]

\subsection*{Theorem (Finite cohomology classifier)}
Let $C^0 \xrightarrow{\delta^0} C^1 \xrightarrow{\delta^1} C^2$ be a finite cochain complex over $\mathbb{Q}$. For $r \in C^1$:
\[
[r] \neq 0 \in H^1
\quad\Longleftrightarrow\quad
\delta^1 r = 0
\;\text{ and }\;
r \notin \operatorname{im}\delta^0.
\]

This is the definition of cohomology $H^1 = \ker\delta^1 / \operatorname{im}\delta^0$.

\subsection{General finite regional obstruction theorem}

Let $G$ be a finite connected oriented graph (vertices $V$, edges $E$, no
2-faces). The cochain complex is $C^0\xrightarrow{\delta^0}C^1\xrightarrow{0}0$.
Every $r\in C^1$ is a cocycle. The cycle space $Z_1(G;\mathbb{Q})=\ker(\delta^{0,T})$
has dimension $|E|-|V|+1$.

\textbf{Theorem (General cycle-detection).}
For $r\in C^1(G;\mathbb{Q})$:
\[
[r]=0\in H^1(G;\mathbb{Q})
\;\Longleftrightarrow\;
\langle z,r\rangle=0 \text{ for every } z\in Z_1(G;\mathbb{Q}).
\]

\textit{Proof.}
$(\Rightarrow)$ If $r=\delta^0 b$ then $\langle z,r\rangle=\langle\delta^{0,T}z,b\rangle=0$.
$(\Leftarrow)$ By rank-nullity, $\dim Z_1+\dim\operatorname{im}\delta^0=\dim C^1$.
Since $\operatorname{im}\delta^0\subseteq Z_1^\perp$ and both have the same dimension,
$\operatorname{im}\delta^0=Z_1^\perp$. So $r\perp Z_1$ implies $r\in\operatorname{im}\delta^0$.
$\square$

\textbf{Consequence.}
A residue is non-removable if and only if it has nonzero circulation around at
least one cycle. Equivalently, $\dim H^1(G;\mathbb{Q})=|E|-|V|+1$ equals the
circuit rank of $G$.

Verified on path $P_3$, triangles $C_3$, cycles $C_4$ (our object), $C_5$,
diamond, filled triangle, and $K_4$ by \texttt{general\_obstruction\_classifier.py}.

\subsection{Complete four-cycle classification theorem}

\textbf{Theorem (Complete four-cycle $H^1$ classifier).}
For the four-region loop nerve $N$, the \emph{circulation functional}
\[
q(a,b,c,d) = -a-b-c+d = \langle z, r\rangle, \quad z=(-1,-1,-1,1),
\]
is the complete invariant of the cohomology class: for any $r=(a,b,c,d)\in C^1$,
\[
[r]=0\in H^1(N;\mathbb{Q})
\;\Longleftrightarrow\;
-a-b-c+d=0.
\]

\textbf{Proof.}
$(\Rightarrow)$ If $r=\delta^0 b$ then $\langle z,r\rangle=\langle\delta^{0,T}z,b\rangle=0$, so $q(r)=0$.

$(\Leftarrow)$ If $q(r)=0$, i.e.\ $d=a+b+c$, set
$b^*=(0,\,a,\,a+b,\,a+b+c)$.
Then $\delta^0 b^* = (a,b,c,d)=r$. $\square$

\textbf{Corollary (Actual object).}
$q(1,1,1,-2)=-1-1-1+(-2)=-5\neq 0$, so $[r]\neq 0\in H^1(N;\mathbb{Q})$. $\square$

The theorem classifies every residue on the four-cycle completely.
The actual object is a corollary, not merely an example.

\subsection{Integral cohomology and modular sensitivity}

\textbf{Theorem (Integral period).}
The circulation map $q:\mathbb{Z}^4\to\mathbb{Z}$ induces an isomorphism
$H^1(N;\mathbb{Z})\cong\mathbb{Z}$.
The actual residue satisfies $[r]=-5\in\mathbb{Z}\cong H^1(N;\mathbb{Z})$.

\textbf{Proof.}
$q$ is surjective over $\mathbb{Z}$ since $q(0,0,0,1)=1$.
By the constructive direction of the classifier theorem (which works integrally:
if $d=a+b+c$ then $b^*\in\mathbb{Z}^4$), $\ker q=\operatorname{im}\delta^0$ over
$\mathbb{Z}$. Therefore $\mathbb{Z}^4/\operatorname{im}\delta^0=\mathbb{Z}^4/\ker
q\cong\mathbb{Z}$, and the class of $r$ is $q(r)=-5$. $\square$

\textbf{Theorem (Modular sensitivity).}
For any field $k$:
\[
[r]\neq 0\in H^1(N;k)
\;\Longleftrightarrow\;
\operatorname{char}(k)\nmid 5.
\]

\textbf{Proof.}
Over $k$, the same analysis gives $H^1(N;k)\cong k$ via $q_k$. Then
$[r]=-5\cdot 1_k\in k$, which is zero iff $\operatorname{char}(k)\mid 5$, i.e.\
$\operatorname{char}(k)=5$.

Explicit coboundary certificate over $\mathbb{F}_5$: the vector $b=(0,1,2,3)$
satisfies $\delta^0 b=(1,1,1,3)\equiv(1,1,1,-2)\pmod{5}=r$, so $r$ is a
coboundary over $\mathbb{F}_5$. $\square$

Modular sensitivity is verified computationally for primes $2,3,5,7,11,13,17,19,23$
by \texttt{classification\_theorem.py}, which emits a machine-readable certificate.

\subsection{Base obstruction}
\textbf{Proposition.} $0 \neq [r] \in H^1(N;\mathbb{Q})$.

\textbf{Proof.}
Since $C^2 = 0$, we have $\delta^1 = 0$, so $\delta^1 r = 0$ trivially. It remains to show $r \notin \operatorname{im}\delta^0$.

The coboundary matrix, with columns indexed by regions and rows by edges, is:
\[
\delta^0 =
\begin{pmatrix}
-1 & 1 & 0 & 0 \\
0 & -1 & 1 & 0 \\
0 & 0 & -1 & 1 \\
-1 & 0 & 0 & 1
\end{pmatrix}.
\]

That is, $(\delta^0 b)_e = b_j - b_i$ for oriented edge $U_iU_j$.

Suppose for contradiction that $\delta^0 b = r$ for some $b = (b_1, b_2, b_3, b_4)$. Then:
\[
b_2 - b_1 = 1, \qquad b_3 - b_2 = 1, \qquad b_4 - b_3 = 1.
\]

Adding: $b_4 - b_1 = 3$.

But the fourth row requires $b_4 - b_1 = -2$. Contradiction. $\square$

\subsection{Cycle witness (independent positive proof)}
The same fact follows from the cycle-pairing duality without solving any linear system.

The vector $z = (-1,-1,-1,1)$ satisfies $\delta^{0,T} z = 0$, i.e., it is a cycle ($\partial z = 0$). Computing:
\[
\langle z, r \rangle
= (-1)(1) + (-1)(1) + (-1)(1) + (1)(-2)
= -5.
\]

If $r = \delta^0 b$ were a coboundary, then:
\[
\langle z, r \rangle = \langle z, \delta^0 b \rangle = \langle \delta^{0,T} z, b \rangle = \langle 0, b \rangle = 0.
\]

Since $\langle z, r \rangle = -5 \neq 0$, we conclude $r \notin \operatorname{im}\delta^0$. $\square$

\subsection{Presentation invariance}
The verdict $[r] \neq 0$ is stable under the following presentation changes:
\begin{itemize}
\item Region renaming (permutation on $C^0$ and $C^1$).
\item Edge reordering (permutation on $C^1$).
\item Edge orientation reversal (diagonal $\pm1$ matrix on $C^1$).
\item Nonzero rational scaling ($r\mapsto\lambda r$, $\lambda\in\mathbb{Q}^\times$).
\item Gauge perturbation ($r\mapsto r+\delta^0 b$).
\end{itemize}

Each is an isomorphism on the cochain complex (or preserves cohomology class), so nonzero classes remain nonzero.

\subsection{Refinement persistence (proved for declared refinements)}
Let $\rho^*:C^1(N)\to C^1(N')$ be a transfer map and $r' = \rho^* r$. If there exists a cycle $z' \in C_1(N')$ with $\partial z' = 0$ and $\langle z', r' \rangle \neq 0$, then $[r'] \neq 0 \in H^1(N')$ (proof by pairing with boundary).

\subsection{Cycle-Lift Persistence Theorem}
Suppose $z \in Z_1(N;\mathbb{Q})$ with $\langle z, r \rangle \neq 0$. If there exists $z'\in Z_1(N';\mathbb{Q})$ and $\lambda\in\mathbb{Q}^\times$ with $\rho_* z' = \lambda z$, then $[\rho^* r]\neq0$ in $H^1(N';\mathbb{Q})$.

\textbf{Proof.} If $\rho^* r = \delta'^0 b'$ then $\langle z', \rho^* r\rangle = \langle \partial' z', b'\rangle = 0$, contradicting adjointness and $\rho_* z' = \lambda z$.

\subsection{Rank criterion and classification table}
Let $B'$ be the boundary matrix of $N'$, $P=\rho_*$, and $K$ a basis for $\ker B' = Z_1(N')$. Then $\rho$ is cycle-faithful relative to $z$ iff $\operatorname{rank}(PK) = \operatorname{rank}([PK\; z])$.

Further details and explicit computations for the four declared refinements are given in the repository and the main manuscript.

\subsection{Residue-Admissibility Bridge Theorem}

\textbf{Setup.}
A \emph{global admissible claim} $\Phi$ is an assignment $\Phi_i\in\mathbb{Q}$
to each region, consistent with all seam data: $\Phi_j-\Phi_i=r_{ij}$ at every
overlap. In cochain terms, $\delta^0\Phi=r$.

\textbf{Theorem (Residue-Admissibility Bridge).}
For $r\in C^1(N;\mathbb{Q})$ satisfying $\delta^1 r=0$:
\begin{enumerate}
\item Gauge-admissible: $\exists\,b\in C^0$ with $\delta^0 b=r$.
\item Globally consistent: $\exists\,\Phi$ with $\Phi_j-\Phi_i=r_{ij}$.
\item Cohomologically trivial: $[r]=0\in H^1(N;\mathbb{Q})$.
\end{enumerate}
are equivalent. When $[r]\neq0$, no global consistent claim exists.
The class $[r]$ is the \emph{warrant debt}: the quantified, irremovable
inconsistency.

\textbf{Proof.} (1)$\Leftrightarrow$(3): definition of $H^1$.
(2)$\Leftrightarrow$(1): $\delta^0\Phi=r\Leftrightarrow\Phi$ is a gauge.
$\square$

\textbf{Warrant Debt Decomposition.}
The orthogonal splitting $C^1=\operatorname{im}\delta^0\oplus Z_1$ yields
$r=r^{\mathrm{adm}}+r^{\mathrm{debt}}$:
\[
r^{\mathrm{adm}}\in\operatorname{im}\delta^0
\;\text{(closest consistent residue)},
\quad
r^{\mathrm{debt}}\in Z_1
\;\text{(irremovable warrant debt)}.
\]
For $r=(1,1,1,-2)$: $r^{\mathrm{adm}}=(-\tfrac{1}{4},-\tfrac{1}{4},-\tfrac{1}{4},-\tfrac{3}{4})$,
$r^{\mathrm{debt}}=(\tfrac{5}{4},\tfrac{5}{4},\tfrac{5}{4},-\tfrac{5}{4})=-\tfrac{5}{4}\,z$,
$\|r^{\mathrm{debt}}\|^2=\tfrac{25}{4}$.

Verified by \texttt{admissibility\_bridge.py}.
